\begin{ExerciseBlock}
\ExerciseHeading{04}{Inelastischer Drehstoß}{angular_momentum_inelastic_angular_collision}
Betrachten Sie folgenden „Drehstoß“: Zwei Scheiben mit Trägheitsmomenten $J_1 = 0,10 \,$kg m² und $J_2 = 0,20 \,$kg m² rotieren anfänglich reibungsfrei auf einer gemeinsamen
Welle gegensinnig mit den Drehzahl (Frequenz) $f_1 = 2\,$1/s und $f_2 = 2,5\,$1/s. Dann werden sie gegeneinander gedrückt und dadurch gekuppelt. Durch die Reibung zwischen den
Kupplungsbelägen der Scheiben stellt sich schließlich eine gemeinsame Drehzahl $f$ ein!

\begin{enumerate}
	\item Wie groß ist die gemeinsame Drehzahl $f$?

	\item Wie viel Energie $\Delta Q$ wird bei diesem Drehstoß in Wärme umgesetzt?
\end{enumerate}
\end{ExerciseBlock}
\UseExerciseSolution{angular_momentum_inelastic_angular_collision}
